\section{Conclusion and Future Work}
\label{sec:conclusion}

In this study, I introduced a highly efficient, targeted approach to detecting deepfakes. By moving away from the computationally expensive trend of massive RGB semantic analysis, I showed that examining the underlying digital physics of an image is an incredibly viable alternative.

My pipeline successfully forces neural networks to act as digital forensic analysts by isolating JPEG compression inconsistencies through Error Level Analysis. That being said, teaching a modern CNN to interpret this highly abstract data was far from simple. My early setbacks highlighted the severe domain shift between standard natural imagery and abstract compression grids. To overcome this hurdle, I had to fully fine-tune the network and carefully apply differential learning rates. This strategy ensured that my EfficientNet-B0 architecture kept its fundamental edge-detection skills while simultaneously adapting to high-frequency noise.

Ultimately, my final results validate my core hypothesis. I successfully built a lightweight model capable of differentiating authentic media from manipulated content, achieving an impressive 86.40\% accuracy rate on the FaceForensics++ dataset.
